\subsection{根系的室与基}

对所有 \(\alpha \in \mathbb{R}\)，令 \(L_{\alpha}\) 为 \(V\) 中由在 \(s_{\alpha}\) 下不变的点组成的超平面。由 \(V\) 中的这些 \(L_{\alpha}\) 所决定的室（Chap. V, § 1, no. 3）称为 \(R\) 的室。由 \(V \to V^{*}\) 这一由内积 \((x|y)\) 定义的双射将 \(\alpha\) 映为 \(\frac{2\alpha^{-}}{(\alpha^{-}|\alpha^{-})}\) 对 \(\alpha \in R\) 成立，因而将 \(L_{\alpha}\) 映为 \(L_{\alpha^{-}}\)，并因而将 \(R\) 的室映为 \(R^{-}\) 的室。若 \(C\) 是 \(R\) 的室，对应的 \(\mathbb{R}^{\sim}\) 的室记为 \(\mathbb{C}^{\sim}\)。由第 2 号的 Prop. 7，\(\mathbb{C}^{\sim}\) 只取决于 \(\mathbb{C}\)，而不取决于 \((x|y)\) 的选择。

定理 2. (i) 群 W(R) 在室的集合上单纯传递地作用。

\begin{enumerate}
\def\labelenumi{(\roman{enumi})}
\setcounter{enumi}{1}
\item
  设 \(\mathrm{C}\) 是一个室。于是 \(\overline{\mathrm{C}}\) 是 \(\mathrm{W}(\mathrm{R})\) 的基本域。
\item
  C 是一个开单纯形锥（Chap. V, § 1, no. 6）。
\item
  设 \(L_{1}, L_{2}, \ldots, L_{l}\) 是 C 的壁。对所有 i，存在唯一的不可分根 \(\alpha_{i}\)，使得 \(L_{i} = L_{\alpha_{i}}\)，并且 \(\alpha_{i}\) 与 C 位于 \(L_{i}\) 的同一侧。
\item
  集合 B(C) = \(\{\alpha_{1},\ldots,\alpha_{l}\}\) 是 V 的一个基。
\item
  C 是由 \(x \in V\) 构成的集合，满足 \(\langle \alpha_i, x \rangle > 0\) 对所有 \(i\) 成立（或等价地，是由 \(x \in V\) 构成的集合，满足 \((x|\alpha_i) > 0\) 对所有 \(i\) 成立）。
\item
  令 \(S\) 为 \(s_{\alpha_i}\) 的集合。则 \((W(R), S)\) 是一个 Coxeter 系统（Chap. IV, § 1, no. 3）。
\end{enumerate}

断言 (i) 和 (vii) 由 Chap. V, § 3, no. 2 的 Th. 1 得到。断言 (ii) 由 Chap. V, § 3, no. 3 的 Th. 2 得到。断言 (iv) 显然。根 \(\alpha_{i}\) 与 \(L_{i}\) 正交，而 \(\alpha_{i}^{-}\) 等同于 \(2\alpha_{i}/(\alpha_{i}|\alpha_{i})\)。由于 W(R) 是本质的（no. 1, Remark 3），断言 (iii)、(v) 和 (vi) 由 Chap. V, § 3, no. 9 的 Prop. 7 得到。

注记。1）断言 (vii) 特别表明，\(\mathrm{W}(\mathbb{R})\) 由反射 \(s_{\alpha_i}\) 生成。

\begin{enumerate}
\def\labelenumi{\arabic{enumi})}
\setcounter{enumi}{1}
\item
  若 \(x, y \in \mathbb{C}\)，则 \((x|y) > 0\)（Chap. V, § 3, no. 5, Lemma 6），换言之，角 \((\widehat{x,y})\) 是锐角。
\item
  令 \(m(\alpha, \beta)\) 为 \(s_{\alpha} s_{\beta} (\alpha, \beta \in \mathrm{B}(\mathrm{C}))\) 的阶。矩阵 \((m(\alpha, \beta))\) 与 (W, S) 的 Coxeter 矩阵（Chap. IV, § 1, no. 9）等同。若 \(\alpha \neq \beta\)，Chap. V, § 3, no. 4 的 Prop. 3 表明，角 \((\widehat{\alpha, \beta})\) 等于 \(\pi - \frac{\pi}{m(\alpha, \beta)}\)；特别地，这个角要么是钝角，要么等于 \(\pi\)，并且 \((\alpha | \beta) \leqslant 0\)。利用第 3 号的列表，得到 \(m(\alpha, \beta)\) 等于 2、3、4 或 6。
\end{enumerate}

定义 2. R 的子集 B 称为 R 的一个基，如果存在 R 的一个室 C，使得 B = B(C)。若 C 是一个室，则 B(C) 称为由 C 定义的 R 的基。

注记。4）Th. 2 的断言 (vi) 表明，映射 C ↦ B(C) 是从室的集合到基的集合的双射。因此，W(R) 在基的集合上单纯传递地作用。

\begin{enumerate}
\def\labelenumi{\arabic{enumi})}
\setcounter{enumi}{4}
\tightlist
\item
  设 C 是 R 的一个室，B 是相应的基。若 \(\alpha\in B\)，置 \(\varphi(\alpha)=\alpha^{-}\)（当 \(2\alpha\notin R\) 时），置 \(\varphi(\alpha)=\frac{1}{2}\alpha^{-}\)（当 \(2\alpha\in R\) 时）。于是 \(\varphi(B)\) 是 \(R^{-}\) 的基，由 \(C^{-}\) 定义；这是因为 \(C^{-}\) 的壁是 \(L_{\alpha^{-}}\)，对 \(\alpha\in B\) 成立。
\end{enumerate}

定义 3. 设 B 是 R 的一个基。R 的 Cartan 矩阵（相对于 B）是矩阵 \((n(\alpha,\beta))_{\alpha,\beta\in\mathbb{B}}\)。

对所有 \(\alpha \in \mathrm{B}\)，\(n(\alpha, \alpha) = 2\)。对于 \(\alpha, \beta \in \mathrm{B}\)，

\[
n (\alpha , \beta) = 2 \frac {(\alpha | \beta)}{(\beta | \beta)} = - 2 \frac {\| \alpha \|}{\| \beta \|} \cos \frac {\pi}{m (\alpha , \beta)},\tag{11}
\]

其中 \(m(\alpha, \beta)\) 如上所述表示 \(s_{\alpha}s_{\beta}\) 的阶。若 \(\alpha \neq \beta\)，则 \(n(\alpha, \beta) = 0\)、-1、-2 或 -3（cf.~no. 3）。

注记。6）Cartan 矩阵 \((n(\alpha,\beta))\) 不应与 Coxeter 矩阵 \((m(\alpha,\beta))\) 混淆。特别要注意，Cartan 矩阵不一定是对称的。

\begin{enumerate}
\def\labelenumi{\arabic{enumi})}
\setcounter{enumi}{6}
\tightlist
\item
  典范编号。若 B 和 \(\mathbf{B}'\) 是 R 的两个基，则存在唯一的 \(w \in \mathbf{W}\)，使得 \(w(\mathbf{B}) = \mathbf{B}'\)。有
\end{enumerate}

\[
n (w (\alpha), w (\beta)) = n (\alpha , \beta) \text { and } m (w (\alpha), w (\beta)) = m (\alpha , \beta)
\]

对于 \(\alpha, \beta \in \mathbf{B}\)。因此，与 B 相关的 Cartan 矩阵和 Coxeter 矩阵可以通过与 \(\mathbf{B}'\) 相关的矩阵同双射的复合得到。

\[
\alpha \mapsto w (\alpha)
\]

从 B 到 B' 得到。

Cartan 矩阵和 Coxeter 矩阵实际上可以按下述方式典范地定义。令 X 是二元组 \((B, \alpha)\) 的集合，其中 B 是 R 的一个基且 \(\alpha \in B\)。群 W 以显然的方式作用于 X，并且 W 在 X 上的每个轨道与集合 \(\{B\} \times B\) 中的每一个恰好交于一点。若 I 是这些轨道的集合，则每个基 B 都有一个典范编号 \((\alpha_{i})_{i \in I}\)。此外，存在唯一的矩阵 \(N = (n_{ij})\)（相应地，\(M = (m_{ij})\)），其类型为 \(I \times I\)，使得对于任意基 B，与 B 相关的 Cartan 矩阵（相应地，Coxeter 矩阵）都可以通过与 B 的典范编号复合而由 N（相应地，M）得到；这个矩阵称为 R 的典范 Cartan 矩阵（相应地，典范 Coxeter 矩阵）。

命题 15. 设 B 是 R 的一个基，\(\alpha\) 是一个不可分根。存在 \(\beta \in B\) 和 \(w \in W(R)\)，使得 \(\alpha = w(\beta)\)。

设 C 是满足 \(\mathrm{B} = \mathrm{B}(\mathrm{C})\) 的室。超平面 \(L_{\alpha}\) 是 R 的某个室 \(C'\) 的一个壁，并且存在 \(\mathrm{W}(\mathrm{R})\) 中的元素将 \(C'\) 变换为 C。因此可以归结到 \(L_{\alpha}\) 是 C 的一个壁的情形。于是 \(\alpha\) 与 R 中的某个元素 \(\beta\) 成比例。由于 \(\alpha\) 和 \(\beta\) 都不可分，\(\alpha = \pm\beta\)。若 \(\alpha = -\beta\)，则 \(\alpha = s_{\beta}(\beta)\)，故命题得证。

推论。设 \(R_{1}\) 和 \(R_{2}\) 是分别位于向量空间 \(V_{1}\) 和 \(V_{2}\) 中的两个约化根系，设 \(B_{1}\) 和 \(B_{2}\) 是 \(R_{1}\) 和 \(R_{2}\) 的基。设 \(f: B_{1} \to B_{2}\) 是一个将 \(R_{1}\) 的 Cartan 矩阵变为 \(R_{2}\) 的 Cartan 矩阵的双射。则存在

存在一个同构 \(\mathrm{F}: \mathrm{V}_1 \to \mathrm{V}_2\)，将 \(\mathbb{R}_1\) 变为 \(\mathbb{R}_2\)，并将 \(\alpha\) 变为 \(f(\alpha)\)，对所有 \(\alpha \in \mathbb{B}_1\) 成立。

设 F 是从 \(V_{1}\) 到 \(V_{2}\) 的同构，将 \(\alpha\) 变为 \(f(\alpha)\)，对所有 \(\alpha \in B_{1}\) 成立。于是 F 将 \(s_{\alpha}\) 变为 \(s_{f(\alpha)}\)，从而将 \(W(R_{1})\) 变为 \(W(R_{2})\)（Th. 2），进而将 \(R_{1}\) 变为 \(R_{2}\)（命题 15）。

命题 16. 设 B 是 R 的一个基，G 是 A(R) 中由保持 B 不变的元素组成的子群。则 W(R) 是 A(R) 的正规子群，而 A(R) 是 G 与 W(R) 的半直积。

若 \(\alpha\in R\) 且 \(t\in\mathrm{A}(\mathrm{R})\)，则 \(ts_{\alpha}t^{-1}=s_{t(\alpha)}\)；由于 \(\mathrm{W}(\mathrm{R})\) 由 \(s_{\alpha}\) 生成，我们看到 \(\mathrm{W}(\mathrm{R})\) 是 \(\mathrm{A}(\mathrm{R})\) 的正规子群。由结构输运，\(\mathrm{A}(\mathrm{R})\) 将 R 的一个基变换为 R 的一个基。由于 \(\mathrm{W}(\mathrm{R})\) 在基的集合上单纯传递地作用，\(\mathrm{A}(\mathrm{R})\) 的每个元素都可以唯一地写成 \(g_{1}g_{2}\) 的形式，其中 \(g_{1}\in\mathrm{W}(\mathrm{R})\)，\(g_{2}\in G\)。

注记。8）设 \(\mathbb{R}_1, \ldots, \mathbb{R}_p\) 是分别位于向量空间 \(\mathbb{V}_1, \ldots, \mathbb{V}_p\) 中的根系，\(\mathbb{R}\) 是 \(\mathbb{R}_i\) 在 \(\mathbb{V} = \prod_i \mathbb{V}_i\) 中的直和，\(\mathbb{C}_i\) 是 \(\mathbb{R}_i\) 的一个室，且 \(\mathbb{B}_i = \mathbb{B}(\mathbb{C}_i)\)。显然 \(\mathbb{C} = \prod_i \mathbb{C}_i\) 是 \(\mathbb{R}\) 的一个室，且 \(\mathbb{B}(\mathbb{C}) = \bigcup_i \mathbb{B}_i\)。由 Th. 2，\(\mathbb{R}\) 的所有室和基都以这种方式得到。

